\documentclass[11pt,a4paper]{article}
\usepackage[utf8]{inputenc}
\usepackage[T1]{fontenc}
\usepackage[ngerman]{babel}
\usepackage{geometry}
\geometry{margin=2.5cm}
\usepackage{booktabs}
\usepackage{enumitem}
\usepackage{xcolor}
\usepackage{scrextend}
\usepackage{hyperref}
\usepackage{tabularx}
% Für mehrseitige Tabellen (optional, hier nicht zwingend benötigt)
\usepackage{longtable}
\hypersetup{
	colorlinks=true,
	linkcolor=blue,
	citecolor=blue,
	urlcolor=blue,
}

\title{Risikovergleich gängiger Drogen in Deutschland \\ mit Schwerpunkt Jugendschutz}
\author{Zusammenstellung auf Basis aktueller Forschung}
\date{\today}

\begin{document}
	
	\maketitle
	
	\begin{abstract}
		Die Frage nach der Gefährlichkeit von Drogen ist komplex, da verschiedene Substanzen auf unterschiedliche Weise Schaden anrichten können. Aus einer medizinischen und gesellschaftlichen Gesamtperspektive wird Alkohol von der überwiegenden Mehrheit der wissenschaftlichen Studien als \textbf{deutlich gefährlicher als Cannabis} eingestuft. Dieses Dokument erweitert den Vergleich um weitere in Deutschland geläufige Substanzen (Tabak, Opioide, Kokain, Amphetamine, MDMA) und erläutert die zentralen gesetzlichen Regelungen zum Jugendschutz.
	\end{abstract}
	
	\section{Einleitung}
	
	Während Alkohol in vielen Gesellschaften tief verwurzelt ist, gilt Cannabis oft als relativ harmlos. Ein direkter Vergleich der Risiken zeigt jedoch ein klares Bild: Alkohol ist in nahezu jeder Kategorie gefährlicher – für das Individuum und für die Gesellschaft. Neben diesen beiden Substanzen existieren jedoch weitere Drogen mit teils erheblich abweichenden Risikoprofilen. Zudem ist der Schutz von Kindern und Jugendlichen ein zentrales Anliegen der deutschen Drogenpolitik, das durch gestaffelte Altersgrenzen und strafbewehrte Verbote umgesetzt wird. Dieser Text fasst die wissenschaftliche Evidenz zu den Risiken sowie die geltenden Jugendschutzbestimmungen zusammen.
	
	\section{Akute Toxizität und Überdosierung}
	
	\begin{itemize}[leftmargin=*, noitemsep]
		\item \textbf{Alkohol:} Eine akute Alkoholvergiftung kann tödlich enden. Weltweit und auch in Deutschland (ca. 44.000 Todesfälle jährlich\cite{who-deaths}\cite{dhs-deaths}) sind Todesfälle durch Alkoholkonsum eine traurige Realität. Schätzungen der WHO zufolge wird in der Europäischen Region jeder elfte Todesfall durch Alkohol verursacht\cite{who-deaths}.
		\item \textbf{Cannabis:} Eine tödliche Überdosis ist extrem unwahrscheinlich. Das \textit{National Institute on Drug Abuse} (NIDA) bestätigt, dass bisher kein Todesfall ausschließlich auf eine Cannabis-Überdosis zurückzuführen ist\cite{nida-no-overdose}. Auch in Deutschland sind keine entsprechenden Fälle verzeichnet\cite{bmg-cannabis}.
	\end{itemize}
	
	\section{Langfristige körperliche Schäden}
	
	\begin{itemize}[leftmargin=*, noitemsep]
		\item \textbf{Alkohol:} Alkohol ist ein Zellgift und eine Hauptursache für Leberzirrhose, Pankreatitis, Herzmuskelerkrankungen und verschiedene Krebsarten. Die WHO stuft Alkohol als \textbf{Gruppe-1-Karzinogen} ein (höchste Risikogruppe, gemeinsam mit Asbest und Tabak)\cite{who-carcinogen}. Er ist kausal mit mindestens sieben Krebsarten verbunden, darunter Mundhöhlen-, Rachen-, Speiseröhren-, Leber-, Kehlkopf-, Darm- und Brustkrebs\cite{who-carcinogen}.
		\item \textbf{Cannabis:} Die direkten körperlichen Schäden sind weniger schwerwiegend. Das größte Risiko für chronische Schäden besteht durch das \textbf{Rauchen} (oft mit Tabak gemischt), was zu Atemwegserkrankungen führen kann. Es gibt Hinweise auf ein erhöhtes Risiko für Herz-Kreislauf-Erkrankungen\cite{nida-heart}. Ein eindeutiger Nachweis für ein Krebsrisiko durch Cannabis allein liegt nicht vor.
	\end{itemize}
	
	\section{Psychische Gesundheit und Gehirn}
	
	\begin{itemize}[leftmargin=*, noitemsep]
		\item \textbf{Alkohol:} Chronischer Konsum kann zu alkoholbedingter Demenz, Gedächtnisstörungen, Depressionen und Angststörungen führen. Ein körperlicher Alkoholentzug kann \textbf{lebensbedrohlich} sein (Delirium tremens).
		\item \textbf{Cannabis:} Das zentrale Risiko ist ein \textbf{erhöhtes Psychose-Risiko}, besonders bei Jugendlichen und jungen Erwachsenen mit entsprechender Veranlagung. Eine kanadische Studie zeigte, dass Jugendliche ein 11-fach höheres Risiko für eine psychotische Störung haben, wenn sie Cannabis konsumieren\cite{cannabis-psychosis}. Hochpotentes Cannabis verdoppelt das Risiko für psychotische Symptome im Vergleich zu niedrigpotentem Cannabis zusätzlich\cite{cannabis-high-potency}. Zudem kann Cannabis die Gehirnentwicklung bei Jugendlichen beeinträchtigen.
	\end{itemize}
	
	\section{Abhängigkeitspotenzial}
	
	\begin{itemize}[leftmargin=*, noitemsep]
		\item \textbf{Alkohol:} Hat ein hohes körperliches und psychisches Abhängigkeitspotenzial. In Deutschland sind schätzungsweise 1,6 Millionen Menschen alkoholabhängig\cite{dhs-abhaengigkeit}.
		\item \textbf{Cannabis:} Führt vor allem zu einer psychischen Abhängigkeit, körperliche Entzugserscheinungen sind milder. Etwa \textbf{9\% aller Konsumenten} entwickeln eine Abhängigkeit, wobei die Rate bei jugendlichen Konsumenten auf 17\% steigt\cite{dhs-cannabis-abhaengigkeit}.
	\end{itemize}
	
	\section{Gesellschaftliche und indirekte Schäden}
	
	\begin{itemize}[leftmargin=*, noitemsep]
		\item \textbf{Alkohol:} Verursacht immense gesellschaftliche Kosten durch Unfälle (ca. 30\% aller tödlichen Verkehrsunfälle), Gewalt, Kriminalität und Produktivitätsausfall. Die Deutschen Hauptstelle für Suchtfragen (DHS) beziffert die volkswirtschaftlichen Kosten des Alkoholkonsums in Deutschland auf rund \textbf{57 Milliarden Euro pro Jahr}\cite{dhs-costs}. In umfassenden Schadensskalen erreicht Alkohol Spitzenwerte, Cannabis wird deutlich niedriger bewertet.
		\item \textbf{Cannabis:} Die gesellschaftlichen Schäden sind im Vergleich zu Alkohol geringer. Die volkswirtschaftlichen Kosten sind schwer zu beziffern, liegen aber deutlich unter denen von Alkohol.
	\end{itemize}
	
	\section{Zusammenfassung der Risikoprofile: Alkohol vs. Cannabis}
	
	\begin{table}[h!]
		\centering
		\small                     % Optional: Schrift etwas verkleinern für mehr Platz
		\begin{tabularx}{\textwidth}{@{} l X X @{}}
			\toprule
			\textbf{Kategorie} & \textbf{Alkohol} & \textbf{Cannabis} \\ \midrule
			Akute Toxizität  & Hoch, Überdosis kann tödlich sein & Sehr gering, tödliche Überdosis extrem unwahrscheinlich\cite{nida-no-overdose} \\
			Organschäden      & Schwer (Leber, Herz, Gehirn, Pankreas) & Geringer, Hauptrisiko durch Rauchen (Lunge) \\
			Krebsrisiko       & Eindeutig krebserregend (WHO Gruppe 1)\cite{who-carcinogen} & Kein eindeutiger Nachweis, Risiko durch Tabak-Mischkonsum \\
			Psychische Risiken & Demenz, Depression, lebensgefährlicher Entzug & Erhöhtes Psychoserisiko, besonders bei Jugendlichen\cite{cannabis-psychosis} \\
			Abhängigkeit      & Hoch, körperlich und psychisch\cite{dhs-abhaengigkeit} & Mittel, vorwiegend psychisch (9\%-17\%)\cite{dhs-cannabis-abhaengigkeit} \\
			Gesellschaftsschaden & Sehr hoch (57 Mrd. € jährlich in D)\cite{dhs-costs} & Geringer im Vergleich zu Alkohol \\ \bottomrule
		\end{tabularx}
	\end{table}
	
	\section{Weitere in Deutschland verbreitete Substanzen}
	
	Die folgende Tabelle ergänzt den Risikovergleich um einige der wichtigsten weiteren Substanzen, die in Deutschland konsumiert werden. Die Bewertung basiert auf Daten des Europäischen Drogenberichts 2025\cite{euda-2025} sowie auf Einschätzungen der Deutschen Hauptstelle für Suchtfragen\cite{dhs-drugprofiles}.
	
	\begin{table}[h!]
		\centering
		\small
		\renewcommand{\arraystretch}{1.2}
		\begin{tabularx}{\textwidth}{@{} l X X X X X @{}}
			\toprule
			\textbf{Kategorie} & \textbf{Tabak / Nikotin} & \textbf{Opioide (Heroin)} & \textbf{Kokain / Crack} & \textbf{Amphetamine (Crystal Meth)} & \textbf{MDMA (Ecstasy)} \\ 
			\midrule
			Akute Toxizität & Gering (keine akute Überdosis) & Extrem hoch (tödlicher Atemstillstand) & Hoch (Herzinfarkt, Schlaganfall) & Hoch (Hyperthermie, Kreislaufversagen) & Mittel (Hyperthermie, Hyponatriämie) \\
			Organschäden & Schwer (Lunge, Herz-Kreislauf, viele Krebsarten) & Schäden durch Injektion (Infektionen, Abszesse) & Herz-Kreislauf-Schäden, Kardiomyopathie & Neurotoxizität, schwere Zahnschäden & Mögliche Neurotoxizität, Leberbelastung \\
			Krebsrisiko & Höchstes Krebsrisiko (Gruppe 1) & Kein direkt karzinogen & Kein eindeutig karzinogen & Kein direkt karzinogen & Kein Nachweis \\
			Psychische Risiken & Hohe Abhängigkeit, Entzugserscheinungen & Schwerste Abhängigkeit, starke Entzugssyndrome & Psychose, Aggressivität, Depression & Psychose, kognitive Defizite & Angst, Panik, „Midweek-Blues“ \\
			Abhängigkeit & Extrem hoch (körperlich \& psychisch) & Extrem hoch (körperlich dominant) & Sehr hoch, besonders Crack & Sehr hoch (rasche Toleranz) & Gering (keine körperliche Abhängigkeit) \\
			Gesellschaftsschaden & Sehr hoch (ca. 80 Mrd. €/Jahr) & Hoch (Kriminalität, Gesundheitskosten) & Hoch (Beschaffungskriminalität) & Regional sehr hoch (z.B. Sachsen, Bayern) & Gering im Vergleich \\
			\bottomrule
		\end{tabularx}
	\end{table}
	
	\subsection*{Kurze Erläuterungen}
	
	\begin{itemize}[leftmargin=*, noitemsep]
		\item \textbf{Tabak/Nikotin:} Verursacht jährlich über 127.000 Todesfälle in Deutschland und stellt die mit Abstand tödlichste legale Droge dar\cite{dhs-tabak}.
		\item \textbf{Opioide:} Heroin und synthetische Opioide (Fentanyl) sind für den Großteil der drogenbedingten Todesfälle verantwortlich (ca. 80\%)\cite{drugcom-opioid}.
		\item \textbf{Kokain/Crack:} Crack (rauchbares Kokain) besitzt ein extrem hohes Suchtpotenzial und wird in der wissenschaftlichen Literatur häufig als eine der gefährlichsten Substanzen eingestuft\cite{nutt-scale}.
		\item \textbf{Amphetamine:} Methamphetamin („Crystal Meth“) führt zu rascher Abhängigkeit und schweren neurologischen Schäden; die Substanz ist besonders in grenznahen Regionen zu Tschechien verbreitet\cite{dhs-crystal}.
		\item \textbf{MDMA:} Die Gefahr besteht vor allem in falschem Trinkverhalten (zu wenig isotonische Getränke) sowie in verunreinigten Pillen; bei sachgemäßem Konsum ist die akute Toxizität relativ gering\cite{drugcom-mdma}.
	\end{itemize}
	
	\section{Jugendschutz in Deutschland}
	
	Der Gesetzgeber hat für Kinder und Jugendliche gestufte Schutzbestimmungen erlassen, um die gesundheitlichen Risiken und die Entwicklungsschäden durch psychoaktive Substanzen zu minimieren. Die wichtigsten Regelungen im Überblick:
	
	\subsection{Alkohol (\S~9 JuSchG)}
	\begin{itemize}[leftmargin=*, noitemsep]
		\item \textbf{Unter 14 Jahren:} Abgabe und Konsum jeglicher alkoholischer Getränke verboten.
		\item \textbf{14–15 Jahre:} Erlaubt sind Bier, Wein und Sekt ausschließlich in Begleitung einer personensorgeberechtigten Person.
		\item \textbf{16–17 Jahre:} Bier, Wein und Sekt ohne Begleitung erlaubt; branntweinhaltige Getränke (Spirituosen, alkoholhaltige Mixgetränke) bleiben verboten.
		\item \textbf{Ab 18 Jahren:} Keine Beschränkungen.
	\end{itemize}
	
	\subsection{Tabak und nikotinhaltige Erzeugnisse (\S~10 JuSchG)}
	\begin{itemize}[leftmargin=*, noitemsep]
		\item \textbf{Unter 18 Jahren:} Abgabe und Konsum von Tabakwaren, E-Zigaretten, E-Shishas sowie nikotinhaltigen Liquids grundsätzlich verboten. Das Verbot gilt auch für nikotinfreie E-Zigaretten, sofern diese zum Verdampfen von Liquids bestimmt sind.
	\end{itemize}
	
	\subsection{Cannabis (Konsumcannabisgesetz – KCanG)}
	\begin{itemize}[leftmargin=*, noitemsep]
		\item \textbf{Unter 18 Jahren:} Absolutes Cannabisverbot. Erwerb, Besitz und Anbau sind für Minderjährige strafbar. Die Weitergabe an Minderjährige wird mit Freiheitsstrafe nicht unter zwei Jahren bestraft (\S~34 KCanG).
		\item Der Konsum in unmittelbarer Gegenwart von Minderjährigen sowie in einem Umkreis von 100 Metern zu Schulen, Kindertagesstätten und Spielplätzen ist für Erwachsene verboten (\S~5 KCanG).
		\item Verstöße von Jugendlichen werden an das Jugendamt gemeldet und können zur Anordnung von Frühinterventionsmaßnahmen führen.
	\end{itemize}
	
	\subsection{Illegale Drogen (Betäubungsmittelgesetz – BtMG)}
	\begin{itemize}[leftmargin=*, noitemsep]
		\item Substanzen wie Heroin, Kokain, Amphetamine, Methamphetamin und MDMA unterliegen dem BtMG. Erwerb, Besitz und Handel sind für alle Altersgruppen strafbar.
		\item Die Abgabe von Betäubungsmitteln an Personen unter 18 Jahren oder das Bestimmen eines Minderjährigen zum unerlaubten Umgang mit BtM wird als besonders schwerer Fall geahndet (\S~29a Abs.~1 Nr.~1, \S~30 Abs.~1 Nr.~2 BtMG) und mit Freiheitsstrafe nicht unter einem Jahr bzw. nicht unter zwei Jahren bestraft.
	\end{itemize}
	
	\section{Fazit}
	
	Die wissenschaftliche Evidenz zeigt, dass Alkohol in den meisten medizinisch relevanten Kategorien ein deutlich höheres Risikoprofil aufweist als Cannabis. Im erweiterten Vergleich wird jedoch deutlich, dass Substanzen wie Tabak und vor allem Opioide und Methamphetamin in Teilbereichen noch gravierendere Schäden verursachen. Keine der genannten Drogen ist harmlos. Besonders für Kinder und Jugendliche bergen alle psychoaktiven Substanzen erhebliche Entwicklungsrisiken. Die deutschen Jugendschutzgesetze tragen dem Rechnung, indem sie für legale Drogen gestaffelte Altersgrenzen definieren und für illegale Drogen strikte Verbote mit empfindlichen Strafandrohungen bei Abgabe an Minderjährige vorsehen. Der sicherste Konsum ist in jedem Alter der Verzicht.
	
	\newpage
	\begin{thebibliography}{99}
		\bibitem{who-deaths} Weltgesundheitsorganisation (WHO), 2025: \href{https://www.who.int/europe/de/news/item/24-12-2025-1-in-3-injury-deaths-caused-by-alcohol--who-europe-reiterates-the-harms}{WHO/Europa weist erneut auf die schädlichen Folgen hin}.
		\bibitem{dhs-deaths} Deutsche Hauptstelle für Suchtfragen (DHS), \textit{Jahrbuch Sucht 2026}, Kapitel 2.1.
		\bibitem{nida-no-overdose} National Institute on Drug Abuse (NIDA); zitiert nach \href{https://www.dw.com/de/faktencheck-wie-ungef%C3%A4hrlich-ist-cannabis/a-65853237}{Deutsche Welle (DW), Faktencheck: Wie (un)gefährlich ist Cannabis?}
		\bibitem{bmg-cannabis} Bundesministerium für Gesundheit: \href{https://www.bundesgesundheitsministerium.de/themen/cannabis/faq-cannabis/ueberdosierung.html}{Cannabis: Fragen und Antworten zur Überdosierung}.
		\bibitem{who-carcinogen} Weltgesundheitsorganisation (WHO), 2021: \href{https://www.who.int/europe/de/home/20-10-2021-alcohol-is-one-of-the-biggest-risk-factors-for-breast-cancer}{Alkohol ist einer der größten Risikofaktoren für Brustkrebs}; IARC-Klassifizierung von 1988.
		\bibitem{nida-heart} National Institute on Drug Abuse (NIDA): \href{https://nida.nih.gov/publications/research-reports/marijuana/cardiovascular-health-effects}{Cardiovascular Health Effects of Marijuana}.
		\bibitem{cannabis-psychosis} André McDonald et al., 2025: Studie zum 11-fach erhöhten Psychoserisiko bei jugendlichen Cannabiskonsumenten, zitiert nach \href{https://www.drugcom.de/news/jugendliche-die-cannabis-konsumieren-haben-ein-erhoehtes-risiko-fuer-psychose/}{drugcom.de}.
		\bibitem{cannabis-high-potency} Lindsey Hines et al., 2024: Längsschnittstudie zu hochpotentem Cannabis und Psychoserisiko, zitiert nach \href{https://www.drugcom.de/news/hochpotenter-cannabis-verdoppelt-risiko-fuer-psychotische-symptome/}{drugcom.de}.
		\bibitem{dhs-abhaengigkeit} Deutsche Hauptstelle für Suchtfragen (DHS): \href{https://www.dhs.de/sucht/zahlen-und-fakten.html}{Zahlen und Fakten zu Alkoholabhängigkeit}.
		\bibitem{dhs-cannabis-abhaengigkeit} Deutsche Hauptstelle für Suchtfragen (DHS), \textit{Jahrbuch Sucht 2025}, Kapitel 2.5; Techniker Krankenkasse: \href{https://www.tk.de/techniker/gesundheit-und-medizin/erste-hilfe-und-gesundheitsvorsorge/sucht-und-abhaengigkeit/cannabis/cannabis-abhaengigkeitspotenzial-2075314}{Cannabis-Abhängigkeitspotenzial}.
		\bibitem{dhs-costs} Deutsche Hauptstelle für Suchtfragen (DHS): \href{https://www.dhs.de/sucht/zahlen-und-fakten.html}{Volkswirtschaftliche Kosten des Alkoholkonsums}.
		% Neue Quellen für erweiterte Substanzen und Jugendschutz
		\bibitem{euda-2025} Europäische Beobachtungsstelle für Drogen und Drogensucht (EUDA), \textit{Europäischer Drogenbericht 2025}.
		\bibitem{dhs-drugprofiles} Deutsche Hauptstelle für Suchtfragen (DHS): \href{https://www.dhs.de/suechte}{Übersicht zu Stoffgruppen und Abhängigkeitsformen}.
		\bibitem{dhs-tabak} Deutsche Hauptstelle für Suchtfragen (DHS): \href{https://www.dhs.de/suechte/tabak}{Tabak – Zahlen und Fakten}.
		\bibitem{drugcom-opioid} \href{https://www.drugcom.de/drogen/opioide/}{drugcom.de: Opioide}.
		\bibitem{nutt-scale} David Nutt et al., 2010: Drug harms in the UK: a multicriteria decision analysis, \textit{The Lancet}.
		\bibitem{dhs-crystal} Deutsche Hauptstelle für Suchtfragen (DHS): \href{https://www.dhs.de/suechte/crystal-meth}{Crystal Meth – Basisinformationen}.
		\bibitem{drugcom-mdma} \href{https://www.drugcom.de/drogen/ecstasy/}{drugcom.de: Ecstasy (MDMA)}.
	\end{thebibliography}
\end{document}